\section*{Bab \ref{ch:b2:series} --- Barisan dan Deret}

\begin{solution}{exo:b2:series:1}
$\sum\frac{\cos n}{n}$: uji Abel dengan $a_n = \frac1n$ dan $b_n =
\cos n$, yang jumlah parsialnya terbatas (yakni bagian real sebuah
jumlah geometri, seperti pada \cref{ex:b2:series:sinn}): jadi konvergen
(tetapi tidak \omterm{def:b2:series:def}{mutlak}, lewat siasat $\cos^2$ yang sama).

$\sum \frac{(-1)^n}{\ln n}$ ($n \geq 2$): lewat uji berselang-seling,
$\frac{1}{\ln n} \downarrow 0$: jadi konvergen; tetapi tidak \omterm{def:b2:series:def}{mutlak}
(sebab $\ln n \leq n$).

$\sum \frac{(-1)^n}{n^{3/4} + \cos n}$: uraikan,
\[
\frac{(-1)^n}{n^{3/4} + \cos n}
= \frac{(-1)^n}{n^{3/4}}\cdot
\frac{1}{1 + \frac{\cos n}{n^{3/4}}}
= \frac{(-1)^n}{n^{3/4}}
- \frac{(-1)^n\cos n}{n^{3/2}} + O\Bigl(\frac{1}{n^{9/4}}\Bigr).
\]
Deret pertamanya: berselang-seling, jadi konvergen. Deret keduanya:
konvergen \omterm{def:b2:series:def}{mutlak} (pada skala $\frac{1}{n^{3/2}}$). Deret ketiganya:
konvergen \omterm{def:b2:series:def}{mutlak}. Totalnya: konvergen.
\end{solution}

\begin{solution}{exo:b2:series:2}
\omterm{thm:b2:series:fubini}{Hasil kali Cauchy} $\sum_{m\geq0} z^m$ dengan dirinya sendiri (keduanya
konvergen \omterm{def:b2:series:def}{mutlak} untuk $\abs z < 1$): koefisien diagonalnya $c_k =
\sum_{m=0}^{k} 1 = k + 1$, sehingga
\[
\frac{1}{(1-z)^2} = \sum_{k\geq0} (k+1)z^k .
\]
Setelah dikalikan $z$ lalu diindeks ulang: $\sum_{n \geq 1} n z^n =
\frac{z}{(1-z)^2}$.
\end{solution}

\begin{solution}{exo:b2:series:3}
Misalkan $B_n = \sum_{k \leq n} b_k \to B$. \omterm{thm:b2:series:abel}{Penjumlahan Abel} dengan $a_k = k$:
\[
\sum_{k=1}^{n} k\,b_k = n B_n - \sum_{k=1}^{n-1} B_k
\quad\Longrightarrow\quad
\frac1n \sum_{k=1}^{n} k b_k = B_n - \frac{1}{n}\sum_{k=1}^{n-1}
B_k .
\]
Rata-rata Cesàro dari $(B_k)$ yang konvergen menuju limitnya $B$
(jilid Tahun ke-1), jadi ruas kanannya menuju $B - B = 0$.
\end{solution}

\begin{solution}{exo:b2:series:4}
Jika $\theta \in \pi\Z$: semua sukunya nol --- jadi konvergen secara
sepele. Andaikan $\theta \notin \pi\Z$.

$\alpha > 1$: konvergen \omterm{def:b2:series:def}{mutlak} (didominasi oleh $n^{-\alpha}$).

$0 < \alpha \leq 1$: uji Abel berlaku (sebab $a_n = n^{-\alpha}
\downarrow 0$; dan jumlah parsial $\sin n\theta$ terbatas oleh
$\frac{1}{\abs{\sin(\theta/2)}}$ lewat jumlah geometri): jadi konvergen.
Tidak \omterm{def:b2:series:def}{mutlak}: sebab $\abs{\sin n\theta} \geq \sin^2 n\theta = \frac{1 -
\cos 2n\theta}{2}$, dan $\sum \frac{1 - \cos 2n\theta}{2n^\alpha}$
divergen (karena $\sum n^{-\alpha}$ divergen, sedangkan $\sum
\frac{\cos 2n\theta}{n^\alpha}$ konvergen lewat Abel bila $2\theta
\notin 2\pi\Z$; adapun kasus terkecuali $2\theta \in 2\pi\Z$ berarti
$\theta \in \pi\Z$, yang sudah ditangani). Jadi semi-konvergen.
\end{solution}

\begin{solution}{exo:b2:series:5}
Misalkan $p_1, p_2, \dots$ suku tak negatif $(u_n)$ menurut urutannya,
dan $q_1, q_2, \dots$ suku negatifnya. Baik $\sum p_k$ maupun $\sum
q_k$ divergen: sebab seandainya salah satunya konvergen, yang lain akan
sama dengan $\sum u_n$ yang konvergen dikurangi ia, jadi konvergen pula
--- dan ketika itu $\sum \abs{u_n} = \sum p_k - \sum q_k$ akan konvergen,
yang bertentangan dengan hipotesisnya. Selain itu $u_n \to 0$ (sebab
$\sum u_n$ konvergen).

Penataan ulang secara rakus: ambil suku positif $p_1, p_2, \dots$ sampai
total berjalannya pertama kali melampaui $\ell$ (mungkin, sebab $\sum p_k
= +\infty$); lalu suku negatif sampai totalnya pertama kali turun di
bawah $\ell$ (mungkin, sebab $\sum q_k = -\infty$); ulangi selamanya
(tiap tahapnya hingga, dan tiap suku terpakai tepat sekali: jadi sebuah
penataan ulang yang sejati). Setelah tiap pergantian, jarak total
berjalannya ke $\ell$ paling banyak sebesar suku terakhir yang dipakai;
dan karena suku yang dipakai pada pergantian ke-$m$ berindeks
$\to \infty$, sedangkan $u_n \to 0$, total berjalannya konvergen ke $\ell$.
\end{solution}

\begin{solution}{exo:b2:series:6}
Keterjumlahannya: jumlah parsial hingga
$\frac{\abs x^{m+n}}{m!n!}$ terbatas oleh $\bigl(\sum_m
\frac{\abs x^m}{m!}\bigr)^2 = \eu^{2\abs x}$. Pengelompokan diagonal
(\cref{thm:b2:series:fubini}) memberi
\[
(\eu^{x})^2 = \sum_{m,n} \frac{x^{m+n}}{m!\,n!}
= \sum_{k=0}^{\infty} x^k \sum_{m+n=k} \frac{1}{m!\,n!}
= \sum_k \frac{x^k}{k!}\sum_{m=0}^{k}\binom km
= \sum_k \frac{(2x)^k}{k!} = \eu^{2x} .
\]
\end{solution}

\begin{solution}{exo:b2:series:7}
Keterjumlahannya: terbatas oleh $\zeta(2)^2$ sebagai keluarga hasil kali
(lewat hujah hasil kali Cauchy pada \cref{thm:b2:series:fubini}).
Pemetaan $(q, a, b) \mapsto (qa, qb)$, dari tripel dengan $\gcd(a, b) =
1$ ke pasangan $(m, n)$, merupakan bijeksi (ambil $q = \gcd(m,n)$).
Setelah \omterm{def:b2:series:summable}{keluarga terjumlahkan} itu dikelompokkan demikian (yakni sebuah
pemilahan himpunan indeksnya --- yang sah bagi \omterm{def:b2:series:summable}{keluarga terjumlahkan}
menurut \cref{thm:b2:series:rearrangement}/\cref{thm:b2:series:fubini}
yang diterapkan pada pemilahan menjadi \omterm{def:b2:structures:countable}{terbilang} banyak kelas):
\[
\zeta(2)^2 = \sum_{q} \sum_{\gcd(a,b)=1} \frac{1}{q^4 a^2 b^2}
= \zeta(4)\, S,
\qquad\text{jadi}\qquad
S = \frac{\zeta(2)^2}{\zeta(4)} .
\]
(Dengan nilai $\zeta(2) = \frac{\pi^2}{6}$ dan
$\zeta(4) = \frac{\pi^4}{90}$ dari
\cref{ch:b2:fourier}: $S = \frac{5}{2}$.)
\end{solution}

\begin{solution}{exo:b2:series:8}
Misalkan $A = \sum a_n$ (yang \omterm{def:b2:series:def}{mutlak}), $B_m = \sum_{k\leq m} b_k \to B$,
yang terbatas oleh $M$. Maka
\[
C_N = \sum_{k=0}^{N} c_k = \sum_{n=0}^{N} a_n B_{N-n}
\]
(kumpulkan menurut indeks $a$). Tulislah
\[
C_N - AB = \sum_{n=0}^{N} a_n (B_{N-n} - B) - B\sum_{n > N} a_n .
\]
Suku terakhirnya menuju $0$. Pecahlah jumlahnya di $n = \lfloor N/2
\rfloor$: untuk $n \leq N/2$ berlaku $N - n \geq N/2$, jadi
$\abs{B_{N-n} - B} \leq \varepsilon_N := \sup_{m \geq N/2}\abs{B_m - B}
\to 0$, sehingga bagian ini $\leq \varepsilon_N \sum\abs{a_n}$;
sedangkan untuk $n > N/2$ berlaku $\abs{B_{N-n} - B} \leq 2M$, sehingga
bagian ini $\leq 2M \sum_{n > N/2} \abs{a_n} \to 0$. Karenanya $C_N \to AB$.
\end{solution}

\begin{solution}{exo:b2:series:9}
Ambil $u_n = 2^{-n} e_n$ (dengan $e_n$ menyatakan barisan yang bernilai
$1$ tunggal pada kedudukan $n$). Maka $\sum \norm{u_n}_\infty = \sum
2^{-n} < \infty$: jadi konvergen \omterm{def:b2:series:def}{mutlak}. Namun jumlah parsialnya $S_N =
(1, \tfrac12, \dots, 2^{-N}, 0, \dots)$ mestinya konvergen ke
barisan $(2^{-n})_n$, yang \emph{tidak} akhirnya nol:
jadi di luar $E$. Di dalam $E$, barisan $(S_N)$ bersifat Cauchy tanpa
limit (sebab $\norm{S_N - x}_\infty \geq 2^{-N-1}$ tak menolong satu pun
$x$ yang akhirnya nol: untuk sembarang $x \in E$ yang nol melampaui
peringkat $K$, berlaku $\norm{S_N - x} \geq 2^{-K-1}$ bila $N > K$):
jadi deretnya tidak konvergen di $E$. Kelengkapan persis itulah yang
diperlukan \cref{thm:b2:nvs:absoluteconvergence}.
\end{solution}

\begin{solution}{exo:b2:series:10}
Baik $\arctan(n+1) - \arctan n$ maupun
$\arctan\frac{1}{n^2+n+1}$ berada di $\intoo{0}{\frac\pi2}$, dan
rumus penjumlahan tangen memberi
\[
\tan\bigl(\arctan(n{+}1) - \arctan n\bigr)
= \frac{(n+1) - n}{1 + n(n+1)} = \frac{1}{n^2 + n + 1} :
\]
jadi tangennya sama pada interval tempat $\tan$ injektif, sehingga
kesamaannya berlaku. Setelah diteleskopkan,
\[
\sum_{n=1}^{N}\arctan\frac{1}{n^2+n+1}
= \arctan(N{+}1) - \arctan 1 \xrightarrow[N\to\infty]{}
\frac\pi2 - \frac\pi4 = \frac\pi4 .
\]

\end{solution}

\begin{solution}{exo:b2:series:11}
Yang pertama: $\sqrt{n+1} - \sqrt n = \frac{1}{\sqrt{n+1} + \sqrt n}
\sim \frac{1}{2\sqrt n}$, jadi sukunya $\sim
2^{-\alpha}n^{-\alpha/2}$: konvergen bila dan hanya bila $\frac\alpha2 >
1$, yakni $\alpha > 2$. Yang kedua: $1 - \cos\frac1n \sim
\frac{1}{2n^2}$: jadi konvergen. Yang ketiga: $\bigl(1 +
\frac1n\bigr)^n = \eu^{\,n\ln(1 + 1/n)} = \eu^{\,1 - \frac1{2n} +
O(n^{-2})} = \eu\bigl(1 - \frac{1}{2n} + O(n^{-2})\bigr)$, sehingga
\[
\eu - \Bigl(1 + \frac1n\Bigr)^{\!n} \sim \frac{\eu}{2n} :
\]
sukunya positif dan setara dengan kelipatan harmonik: jadi divergen.
\end{solution}

\begin{solution}{exo:b2:series:12}
Berkat kemonotonannya, $n\,a_{2n} \leq a_{n+1} + a_{n+2} + \dots +
a_{2n} = S_{2n} - S_n \to 0$ (lewat kriteria Cauchy bagi deret yang
konvergen). Karenanya $2n\,a_{2n} \to 0$, dan $(2n{+}1)\,
a_{2n+1} \leq (2n{+}1)a_{2n} = \frac{2n+1}{2n}\,(2n\,a_{2n}) \to
0$: jadi kedua barisan bagian $(na_n)$ menuju $0$, sehingga $na_n \to 0$.

\emph{Konversnya gagal:} $a_n = \frac1{n\ln n}$ positif dan
turun dengan $na_n = \frac1{\ln n} \to 0$, namun $\sum a_n$
divergen (yakni perbatasan Bertrand, \cref{pb:b2:series:1}, pertanyaan
18). \emph{Kemonotonannya penting:} ambil $a_n = \frac1n$ bila $n$
berupa pangkat $2$ dan $a_n = 2^{-n}$ bila bukan: maka $\sum a_n \leq
\sum_k 2^{-k} + \sum_n 2^{-n} < \infty$, tetapi $na_n = 1$ sepanjang
pangkat $2$: jadi $na_n \not\to 0$.
\end{solution}

\begin{solution}{pb:b2:series:1}
\textbf{1.} Induksi pada $m$: untuk $m = 0$ kedua ruasnya $1$;
lalu langkahnya mengalikan dengan $\cos\theta + \iu\sin\theta$ dan
memakai rumus penjumlahan $\cos(m\theta + \theta) =
\cos m\theta\cos\theta - \sin m\theta\sin\theta$ serta $\sin(m\theta
+ \theta) = \sin m\theta\cos\theta + \cos m\theta\sin\theta$.
Sebagai gantinya, menguraikannya lewat teorema binomial dengan $m =
2n+1$ lalu mengumpulkan bagian imajinernya (yakni pangkat ganjil
$\iu\sin\theta$, dengan $\iu^{2j+1} = (-1)^j\iu$) memberi
\[
\sin\bigl((2n{+}1)\theta\bigr) = \sum_{j=0}^{n}(-1)^j
\binom{2n+1}{2j+1}\cos^{2(n-j)}\theta\,\sin^{2j+1}\theta .
\]

\textbf{2.} Pada $\intoo0{\frac\pi2}$ berlaku $\sin\theta \neq 0$:
keluarkan faktor $\sin^{2n+1}\theta$ dari tiap sukunya, sehingga tersisa
$\bigl(\frac{\cos^2\theta}{\sin^2\theta}\bigr)^{n-j} =
(\cot^2\theta)^{n-j}$: yakni kesamaan yang tertulis tadi dengan $P_n(x) =
\sum_j(-1)^j\binom{2n+1}{2j+1}x^{n-j}$. Koefisien berderajat $n$-nya
berasal dari $j = 0$, yaitu $\binom{2n+1}{1} = 2n + 1 \neq 0$.

\textbf{3.} Di $\theta_k = \frac{k\pi}{2n+1}$:
$\sin\bigl((2n{+}1)\theta_k\bigr) = \sin k\pi = 0$ sedangkan
$\sin^{2n+1}\theta_k \neq 0$, jadi $P_n(\cot^2\theta_k) = 0$. Nilai
$\theta_k$ naik tegas di $\intoo0{\frac\pi2}$, tempat
$\cot^2$ turun tegas: jadi nilai $x_k =
\cot^2\theta_k$ berbeda sepasang demi sepasang --- yakni $n$ akar
berbeda bagi polinomial berderajat $n$, sehingga itulah seluruh akarnya.

\textbf{4.} Menurut Vieta, jumlah akarnya adalah negatif rasio
koefisien $x^{n-1}$ terhadap koefisien $x^n$:
\[
\sum_{k=1}^{n}\cot^2\theta_k =
\frac{\binom{2n+1}{3}}{\binom{2n+1}{1}}
= \frac{(2n+1)(2n)(2n-1)/6}{2n+1} = \frac{n(2n-1)}{3}.
\]

\textbf{5.} Karena $\frac{1}{\sin^2\theta} = 1 + \cot^2\theta$, setelah
dijumlahkan diperoleh $n + \frac{n(2n-1)}3 = \frac{3n + 2n^2 - n}{3} =
\frac{2n(n+1)}{3}$.

\textbf{6.} Pada $\intoo{0}{\frac\pi2}$: $\sin\theta < \theta <
\tan\theta$ (jilid Tahun ke-1). Mengambil kebalikannya membalik
urutannya: $\cot\theta < \frac1\theta < \frac1{\sin\theta}$, lalu
mengkuadratkannya (semuanya positif) memberi $\cot^2\theta <
\frac1{\theta^2} < \frac1{\sin^2\theta}$.

\textbf{7.} Jumlahkan pertanyaan 6 di $\theta = \theta_k$ atas $k \leq
n$, dengan memakai pertanyaan 4 dan 5, serta $\frac1{\theta_k^2} =
\frac{(2n+1)^2}{k^2\pi^2}$:
\[
\frac{n(2n-1)}3 < \frac{(2n+1)^2}{\pi^2}\sum_{k=1}^n\frac1{k^2}
< \frac{2n(n+1)}3 .
\]

\textbf{8.} Kalikan dengan $\frac{\pi^2}{(2n+1)^2}$:
\[
\frac{\pi^2}{3}\cdot\frac{n(2n-1)}{(2n+1)^2}
< \sum_{k=1}^{n}\frac1{k^2} <
\frac{\pi^2}{3}\cdot\frac{2n(n+1)}{(2n+1)^2}.
\]
Kedua batasnya menuju $\frac{\pi^2}3\cdot\frac12 = \frac{\pi^2}6$
(sebab pecahan rasionalnya menuju $\frac12$). Jumlah parsialnya
naik, jadi ia konvergen, dan apitannya memberi $\zeta(2) =
\frac{\pi^2}6$: itulah teorema Euler, lewat bukti Cauchy.

\textbf{9.} Jumlah parsialnya naik ke $\zeta(2) = \frac{\pi^2}6$, jadi
$0 \leq \frac{\pi^2}6 - \sum_{k\leq n}k^{-2}$;
lalu batas bawah pada pertanyaan 8 memberi
\[
\frac{\pi^2}6 - \sum_{k\leq n}\frac1{k^2}
\leq \frac{\pi^2}6 - \frac{\pi^2}3\cdot\frac{n(2n-1)}{(2n+1)^2}
= \frac{\pi^2}6\cdot\frac{(2n+1)^2 - (4n^2 -
2n)}{(2n+1)^2}
= \frac{\pi^2}6\cdot\frac{6n + 1}{(2n+1)^2}
= O\Bigl(\frac1n\Bigr),
\]
yang cocok dengan ekor eksak $\frac1n - \frac1{2n^2} + O(n^{-3})$ pada
\cref{exo:b2:comparison:11}.

\textbf{10.} Fungsi simetrik elementer kedua atas akarnya adalah
$\sigma_2 = \frac{\binom{2n+1}5}{\binom{2n+1}1} =
\frac{(2n)(2n-1)(2n-2)(2n-3)}{120}$, sehingga
\[
\sum_k\cot^4\theta_k = \sigma_1^2 - 2\sigma_2
= \Bigl(\frac{n(2n-1)}3\Bigr)^2 -
\frac{2n(2n-1)(2n-2)(2n-3)}{60}
\sim \frac{4n^4}9 - \frac{4n^4}{15} = \frac{8n^4}{45}.
\]
Setelah $\cot^4\theta < \theta^{-4} < (1 + \cot^2\theta)^2 =
1 + 2\cot^2\theta + \cot^4\theta$ diapit lalu dijumlahkan: kedua jumlah
luarnya bernilai $\frac{8n^4}{45}(1 + o(1))$ (sebab tambahan $n +
2\sigma_1 = O(n^2)$ terabaikan), sedangkan jumlah tengahnya adalah
$\frac{(2n+1)^4}{\pi^4}\sum_{k\leq n}k^{-4}$. Karenanya
\[
\sum_{k\leq n}\frac1{k^4} \longrightarrow
\pi^4\cdot\frac{8/45}{16} = \frac{\pi^4}{90}.
\]

\textbf{11.} Setelah $\zeta(2)$ dipecah menurut paritasnya:
$\sum_{\text{genap}} = \sum_j\frac{1}{(2j)^2} = \frac14\zeta(2) =
\frac{\pi^2}{24}$, jadi $\sum_{\text{ganjil}} = \zeta(2) -
\frac{\pi^2}{24} = \frac{\pi^2}8$. Untuk yang berselang-seling:
$\sum_k\frac{(-1)^{k-1}}{k^2} = \sum_{\text{ganjil}} -
\sum_{\text{genap}} = \frac{\pi^2}8 - \frac{\pi^2}{24} =
\frac{\pi^2}{12}$ (dan kekonvergenan \omterm{def:b2:series:def}{mutlaknya} mengesahkan
pengelompokan ulang itu, \cref{thm:b2:series:rearrangement}).

\textbf{12.} $S = \frac{\zeta(2)^2}{\zeta(4)} =
\frac{(\pi^2/6)^2}{\pi^4/90} = \frac{90}{36} = \frac52$.
Heuristiknya: kesamaan $\zeta(2)^2 = \zeta(4)S$ pada
\cref{exo:b2:series:7} mengatakan bahwa mengeluarkan $\gcd$
menormalkan ulang pasangannya menjadi pasangan yang saling prima;
lalu kebalikannya $\frac1{\zeta(2)} = \frac6{\pi^2} \approx 0.608$
menjadi calon alami bagi kerapatan pasangan saling prima di antara
semua pasangan --- yakni pernyataan tentang $\lim_N
\frac{1}{N^2}\#\{(m,n) \leq N : \gcd = 1\}$ yang bukti jujurnya
(\omterm{def:b2:metric:complete}{lengkap} dengan suku galatnya) menjadi bagian jilid Tahun ke-3.

\textbf{13.} Penjumlahan langsung bergalat $\sim \frac1n$: jadi enam
angka menuntut sekitar $10^6$ suku. Sedangkan jumlah terkoreksi
$\sum_{k\leq n}k^{-2} + \frac1n - \frac1{2n^2}$ bergalat
$O(n^{-3})$: di $n = 100$ ia sama dengan $1.6449339\dots$ terhadap
$\frac{\pi^2}6 = 1.6449341\dots$ --- galatnya $1.7\cdot10^{-7}$,
yakni tujuh angka dari seratus suku. Jadi koreksi asimtotik mengalahkan
kesabaran mentah sejauh empat orde besaran.

\textbf{14.} Untuk $n = 1$: $P_1(x) = 3x - 1$, dengan akar $\frac13$, dan
memang $\cot^2\frac\pi3 = \bigl(\frac1{\sqrt3}\bigr)^2 =
\frac13 = \frac{1\cdot1}3$. Untuk $n = 2$: rumusnya meramalkan
$\frac{2\cdot3}3 = 2$; lalu dengan $\cos\frac\pi5 = \frac{1 +
\sqrt5}{4}$ terhitung $\cot^2 36^\circ \approx 1.894$ dan
$\cot^2 72^\circ \approx 0.106$: jumlahnya $2.000$.

\textbf{15.} Bilangan $\tan^2\theta_k = \frac1{x_k}$ merupakan
akar $Q(x) = x^nP_n\bigl(\frac1x\bigr) =
\sum_{j=0}^n(-1)^j\binom{2n+1}{2j+1}x^j$ (sebab $x_k$ tak nol).
Vieta pada $Q$: koefisien utamanya $(-1)^n$
(dari suku $j = n$), koefisien berikutnya $(-1)^{n-1}\binom{2n+1}{2n-1} =
(-1)^{n-1}\binom{2n+1}{2}$, sehingga
\[
\sum_{k=1}^n\tan^2\theta_k =
-\frac{(-1)^{n-1}\binom{2n+1}2}{(-1)^n} = \binom{2n+1}2\cdot
\frac{2}{2n+1}\cdot\frac{2n+1}{2}
= n(2n+1).
\]
Periksa $n = 1$: $\tan^2\frac\pi3 = 3 = 1\cdot3$.

\textbf{16.} Misalkan $\gamma = \frac{1+\alpha}2 \in
\intoo{1}{\alpha}$. Maka $\frac{n^{-\alpha}(\ln
n)^{-\beta}}{n^{-\gamma}} = n^{\gamma - \alpha}(\ln n)^{-\beta}
\to 0$ (sebab pangkat negatif $n$ mengalahkan pangkat berapa pun dari
$\ln n$), jadi akhirnya sukunya $\leq n^{-\gamma}$ dengan $\gamma > 1$:
sehingga konvergen lewat pembandingan dengan deret Riemann.

\textbf{17.} Misalkan $\gamma = \frac{1+\alpha}2 \in
\intoo{\alpha}{1}$: kini $\frac{n^{-\gamma}}{n^{-\alpha}(\ln
n)^{-\beta}} = n^{\alpha-\gamma}(\ln n)^{\beta} \to 0$, jadi
akhirnya sukunya $\geq n^{-\gamma}$ dengan $\gamma < 1$:
sehingga divergen.

\textbf{18.} Fungsi $f(t) = \frac1{t(\ln t)^\beta}$ bersifat positif,
\omterm{def:b2:metric:continuity}{kontinu}, dan turun untuk $t$ yang besar (sebab logaritmanya
berturunan $-\frac1t\bigl(1 + \frac{\beta}{\ln t}\bigr) < 0$
pada akhirnya). Antiturunannya: untuk $\beta \neq 1$,
$\int^x f = \frac{(\ln x)^{1-\beta}}{1-\beta} + \text{tetapan}$,
yang punya limit hingga bila dan hanya bila $\beta > 1$; sedangkan untuk
$\beta = 1$, $\int^x f = \ln\ln x \to \infty$. Menurut
\cref{thm:b2:comparison:seriesintegral}, deret dan
integralnya bersifat sama: jadi konvergen bila dan hanya bila $\beta > 1$.

\textbf{19.} Uji yang sama: $\frac{\dd}{\dd t}\ln\ln\ln t =
\frac{1}{t\ln t\,\ln\ln t}$, dan $\ln\ln\ln t \to \infty$:
jadi divergen. Sedangkan $\frac{\dd}{\dd
t}\Bigl(-\frac1{\ln\ln t}\Bigr) = \frac{1}{t\ln t\,(\ln\ln t)^2}$
dengan $-\frac1{\ln\ln t} \to 0$: jadi konvergen.

\textbf{20.} Yang pertama: $n^{1/\ln n} = \eu^{\ln n/\ln n} = \eu$, jadi
sukunya tepat $\frac{1}{\eu\,n}$: yakni kelipatan deret
harmonik, sehingga \emph{divergen} --- sebab eksponen $1 +
\frac1{\ln n}$ merangkak ke $1$ terlalu cepat. Yang kedua:
$n^{1/\ln\ln n} = \eu^{\ln n/\ln\ln n}$, dan akhirnya $\frac{\ln
n}{\ln\ln n} \geq 2\ln\ln n$, jadi $n^{1/\ln\ln n} \geq (\ln n)^2$:
sehingga sukunya $\leq \frac1{n(\ln n)^2}$, yakni deret Bertrand yang
konvergen (pertanyaan 18): jadi \emph{konvergen}. Perbatasannya lewat
tegas di antara kedua eksponen itu.

\textbf{21.} Kita punya $R_n \downarrow 0$ dan
\[
\sqrt{R_n} - \sqrt{R_{n+1}} = \frac{R_n -
R_{n+1}}{\sqrt{R_n} + \sqrt{R_{n+1}}} =
\frac{a_n}{\sqrt{R_n} + \sqrt{R_{n+1}}} \geq
\frac{a_n}{2\sqrt{R_n}},
\]
jadi $\sum_n \frac{a_n}{\sqrt{R_n}} \leq 2\sum_n(\sqrt{R_n} -
\sqrt{R_{n+1}}) = 2\sqrt{R_1} < \infty$ (lewat teleskop). Namun
$\frac{a_n/\sqrt{R_n}}{a_n} = \frac1{\sqrt{R_n}} \to \infty$:
jadi deret yang baru itu konvergen padahal ia tak hingga kali lebih
besar. Tak ada deret konvergen yang paling lambat; jadi uji perbandingan
terhadap keluarga tetap mana pun tak akan pernah \omterm{def:b2:metric:complete}{lengkap}.

\textbf{22.} Untuk $(a_n)$ positif yang turun, kelompokkan sukunya
di antara pangkat $2$ yang berurutan:
\[
2^{k}a_{2^{k+1}} \leq \sum_{n=2^k}^{2^{k+1}-1} a_n \leq
2^ka_{2^k} .
\]
Setelah dijumlahkan atas $k$: jika $\sum 2^ka_{2^k}$ konvergen, maka
jumlah parsial $\sum a_n$ terbatas (jadi konvergen); sebaliknya jika
$\sum a_n$ konvergen, maka $\sum_k 2^{k+1}a_{2^{k+1}} \leq 2\sum_n a_n <
\infty$. Untuk $a_n = \frac1{n(\ln n)^\beta}$: $2^ka_{2^k} =
\frac{1}{(k\ln 2)^\beta}$, dan $\sum k^{-\beta}$ konvergen bila dan
hanya bila $\beta > 1$: jadi perbatasan pertanyaan 18 lagi, tanpa integral.

\textbf{23.} Menurut perbandingan integralnya,
\[
\sum_{n > N}\frac{1}{n(\ln n)^2} \geq
\int_{N+1}^{\infty}\frac{\dd t}{t(\ln t)^2} =
\frac{1}{\ln(N+1)},
\]
yang di $N = 10^6$ bernilai $\approx 0.0724$: jadi setelah sejuta suku
ekornya masih melampaui $0.07$ --- deretnya konvergen, tetapi tak ada
penjumlahan langsung yang akan pernah memperlihatkan jumlahnya.
Bandingkan dengan pertanyaan 13, tempat satu koreksi asimtotik membeli
tujuh angka dari seratus suku: mengetahui \emph{bagaimana} sebuah deret
konvergen lebih berharga daripada mengetahui bahwa ia konvergen.

\textbf{24.} $\sum\frac1{n\ln n}$: divergen ($\alpha = 1$,
$\beta = 1$, pertanyaan 18). $\sum\frac1{n^{1.01}}$: konvergen
(Riemann, $\alpha > 1$). $\sum\frac{(\ln n)^{100}}{n^{1.001}}$:
konvergen ($\alpha = 1.001 > 1$, $\beta = -100$, pertanyaan 16).
$\sum\frac1{n\ln n(\ln\ln n)^3}$: konvergen (dengan pola pertanyaan
19: antiturunannya $-\frac12(\ln\ln t)^{-2}$, yang berlimit
hingga).

\textbf{25.} De Moivre mengubah penolkan $\sin(2n{+}1)
\theta_k$ menjadi penolkan sebuah polinomial di
$\cot^2\theta_k$, lalu Vieta membaca jumlah akar eksak yang hanya
dapat ditaksir oleh analisis semata (pertanyaan 1--5).
Apitannya memerlukan nilai \emph{eksak} $\frac{n(2n-1)}3$ dan
$\frac{2n(n+1)}3$ pada kedua sisinya --- sebab yang setara justru akan
mengemis pada pertanyaannya, karena seluruh intinya adalah konstanta
$\frac{\pi^2}6$ (pertanyaan 7--8). Bagian IV seluruhnya berjalan pada
perbandingan deret dengan integral pada \cref{ch:b2:comparison}, dengan
antiturunan logaritmiknya yang menggolongkan (pertanyaan
18--19). Pertanyaan 21 menghancurkan impian akan uji
perbandingan semesta: sebab di bawah setiap deret konvergen ada deret
lain yang tak hingga kali lebih lambat --- skala seperti milik Bertrand
memetakan perbatasannya makin halus tetapi tak pernah dapat mencapainya.
Puncaknya: $\zeta(2) = \frac{\pi^2}6$ milik Euler beserta lantai
atasnya $\zeta(4) = \frac{\pi^4}{90}$ (pertanyaan 8 dan 10), dan
penggolongan Bertrand (pertanyaan 16--18); adapun $\zeta(2)$ kembali
pada \cref{ch:b2:fourier}, tempat kesamaan Parseval membuktikannya
ulang dalam satu baris dari deret Fourier gelombang gergaji --- satu
konstanta, dua peradaban.

\end{solution}
